\chapter{Algoritmo Genetico Ibrido (HGA)}

Tra gli algoritmi proposti per il progetto, la scelta \`e ricaduta
sull'\textbf{Algoritmo Genetico Ibrido (HGA)} \cite{holland1975adaptation, goldberg1989genetic}, ispirato al framework HGSADC (Hybrid Genetic Search with Adaptive Diversity Control) \cite{vidal2012unified}.
Gli algoritmi genetici esplorano
efficientemente lo spazio delle soluzioni attraverso i meccanismi di selezione,
crossover e mutazione, riducendo il rischio di convergenza prematura verso minimi
locali. L'ibridazione con operatori di ricerca locale (2-opt, Or-opt, Relocate,
Exchange) \cite{lin1973effective, or1976traveling} permette di migliorare le soluzioni prodotte dal processo evolutivo,
combinando l'esplorazione globale con lo sfruttamento locale.

\section{Encoding e Algoritmo di Split}

L'algoritmo utilizza un \textbf{encoding a permutazione} dei clienti (escluso il depot).
Ogni cromosoma \`e una permutazione $\pi = (\pi_1, \pi_2, \dots, \pi_n)$ dei nodi
$1, 2, \dots, n$. La decodifica da permutazione a soluzione (insieme di route) \`e
effettuata tramite l'\textbf{algoritmo di Split} di Prins \cite{prins2004simple}, che trasforma una
permutazione in un insieme ottimale di route che rispettano i vincoli di capacit\`a,
utilizzando programmazione dinamica con complessit\`a $O(n^2)$.

Sia $\pi = (\pi_1, \dots, \pi_n)$ una permutazione di clienti. Definiamo $V[k]$
come il costo minimo per servire i primi $k$ clienti della permutazione:

$$V[i] = \min_{j: \sum_{k=j}^{i} q_{\pi_k} \leq \sigma} \left[ V[j-1] + d_{0,\pi_j} + \sum_{k=j}^{i-1} d_{\pi_k,\pi_{k+1}} + d_{\pi_i,0} \right]$$

Dopo la DP, si ricostruiscono le route tramite backtracking. L'uso dello split come
decodifica \`e fondamentale: invece di codificare esplicitamente le route nel cromosoma,
si lascia che la DP trovi la partizione ottimale, riducendo lo spazio di ricerca e
garantendo l'ottimalit\`a della decodifica data la permutazione.

La scelta di una permutazione senza delimitatori \`e deliberata. Le rappresentazioni
con delimitatori (es.\ inserire zeri tra le route) producono figli non validi durante
il crossover, richiedendo costose riparazioni. La permutazione semplice, invece,
garantisce che qualsiasi cromosoma sia sintatticamente valido e separa elegantemente
il problema di \textbf{sequenziamento} dei clienti (compito del GA) dal problema di
\textbf{partizionamento in rotte} (compito dello Split).

\begin{algorithm}[H]
\footnotesize
\caption{Split Algorithm (Prins, 2004)}
\label{alg:split}
\KwIn{Permutazione $\pi = (\pi_1,\dots,\pi_n)$, matrice distanze $d$, domande $q$, capacit\`a $Q$, depot $0$}
\KwOut{Rotte ottimali $\mathcal{R}$ e costo totale}

\tcp{Pre-calcolo distanze cumulative lungo $\pi$}
$\mathit{cum}[0] \gets 0$\;
\For{$i \gets 1$ \KwTo $n-1$}{
    $\mathit{cum}[i+1] \gets \mathit{cum}[i] + d[\pi_{i},\pi_{i+1}]$\;
}

\tcp{Programmazione Dinamica}
$V[0] \gets 0$\;
$\mathit{pred}[0] \gets -1$\;
\For{$i \gets 1$ \KwTo $n$}{
    $V[i] \gets +\infty$\;
    $\mathit{load} \gets 0$\;
    \For{$j \gets i$ \KwTo $1$}{
        $\mathit{load} \gets \mathit{load} + q[\pi_j]$\;
        \If{$\mathit{load} > Q$}{
            \textbf{break}\;
        }
        $\mathit{routeCost} \gets d[0,\pi_j] + (\mathit{cum}[i] - \mathit{cum}[j]) + d[\pi_i,0]$\;
        \If{$V[j-1] + \mathit{routeCost} < V[i]$}{
            $V[i] \gets V[j-1] + \mathit{routeCost}$\;
            $\mathit{pred}[i] \gets j-1$\;
        }
    }
}

\tcp{Backtracking per ricostruire le route}
$\mathcal{R} \gets \emptyset$\;
$i \gets n$\;
\While{$i > 0$}{
    $j \gets \mathit{pred}[i] + 1$\;
    $\mathcal{R}.\text{append}(\pi[j \dots i])$\;
    $i \gets j-1$\;
}
$\mathcal{R} \gets \text{reverse}(\mathcal{R})$\;  \tcp{Le route sono state costruite in ordine inverso}
\Return{$\mathcal{R},\; V[n]$}\;
\end{algorithm}

\section{Popolazione Iniziale}

La popolazione iniziale ($\mu = 81$ individui) \`e generata combinando:
\begin{enumerate}[label=\textbf{\arabic*.}]
    \item \textbf{Nearest Neighbor Heuristic}: Costruisce una permutazione partendo
          dal depot e selezionando iterativamente il cliente pi\`u vicino non ancora
          visitato.
    \item \textbf{Clarke-Wright Savings Heuristic} \cite{clarke1964scheduling}: Utilizza l'algoritmo dei risparmi
          per costruire route greedy, poi le appiattisce in una permutazione.
    \item \textbf{Permutazioni casuali}: Le restanti permutazioni sono generate in modo
          random per garantire diversit\`a genetica.
\end{enumerate}

Le due soluzioni euristiche forniscono ``punti di partenza'' di buona qualit\`a,
accelerando la convergenza iniziale. Le permutazioni casuali garantiscono
\textbf{diversit\`a genetica}, evitando la convergenza prematura della popolazione
su un ottimo locale. La proporzione, deliberatamente sbilanciata verso il random
($\sim$97.5\% con $\mu=81$), favorisce l'esplorazione rispetto allo sfruttamento
nella fase iniziale.

\section{Operatori Genetici}

\subsection{Selezione a Torneo}

Viene utilizzata la \textbf{selezione a torneo} \cite{goldberg1991comparative} con $k = 4$. Quattro individui sono
scelti casualmente dalla popolazione e il migliore (costo minore) viene selezionato.
Questo metodo offre un buon bilanciamento tra pressione selettiva e diversit\`a.

Rispetto alla roulette wheel (che soffre di dominanza precoce dei migliori) e alla
selezione per rango (pi\`u lenta), il torneo offre un equilibrio regolabile tramite
$k$ tra \textbf{pressione selettiva} ($k$ alto favorisce i migliori) e
\textbf{diversit\`a} ($k$ basso d\`a pi\`u chance agli individui mediocri). Il
valore $k=4$, determinato da Optuna, offre una pressione selettiva moderata-alta,
adatta al paesaggio di fitness multimodale del CVRP. L'implementazione include un
percorso rapido ottimizzato per il caso $k=2$.

\subsection{Crossover: Order Crossover (OX)}

L'Order Crossover \cite{davis1985applying} preserva l'ordine relativo dei geni da entrambi i genitori.
L'Algoritmo~\ref{alg:ox} descrive il procedimento: vengono selezionati due punti di
taglio casuali, il segmento centrale viene copiato dal primo genitore, e le posizioni
rimanenti sono riempite con i geni del secondo genitore in ordine di apparizione,
saltando quelli gi\`a presenti. Il tasso di crossover \`e $p_c = 0.675$.

OX \`e l'operatore ideale per permutazioni perch\'e: (i) preserva l'ordine relativo
di entrambi i genitori, a differenza del PMX che preserva le posizioni assolute,
meno rilevante nel routing; (ii) non produce duplicati --- ogni cliente appare
esattamente una volta, senza bisogno di riparazioni; (iii) trasmette blocchi
contigui di clienti, che spesso corrispondono a sequenze geograficamente vicine
(buone rotte parziali). \`E l'operatore standard nella letteratura VRP basata su
Split \cite{prins2004simple, vidal2012unified}. Il tasso $p_c = 0.675$, inferiore al
default 0.8, riflette la scelta di Optuna: un crossover leggermente meno frequente,
combinato con pi\`u ricerca locale, produce risultati migliori.

\begin{algorithm}[H]
\footnotesize
\caption{Order Crossover (OX) \cite{davis1985applying}}
\label{alg:ox}
\KwIn{Permutazioni genitore $P_1$, $P_2$ di lunghezza $n$}
\KwOut{Permutazione figlio $C$}

\tcp{Seleziona due punti di taglio casuali}
$cx_1 \gets \text{random}(0, n-1)$\;
$cx_2 \gets \text{random}(0, n-1)$\;
\If{$cx_1 > cx_2$}{
    $cx_1, cx_2 \gets cx_2, cx_1$\;
}

\tcp{Copia il segmento centrale da $P_1$}
$C \gets$ array di $n$ elementi non inizializzati\;
\For{$i \gets cx_1$ \KwTo $cx_2$}{
    $C[i] \gets P_1[i]$\;
}

\tcp{Riempi le posizioni rimanenti con i geni di $P_2$ in ordine}
$\mathit{seg} \gets \{C[cx_1], \dots, C[cx_2]\}$\; \tcp{Insieme dei geni gi\`a presenti}
$j \gets (cx_2 + 1) \bmod n$\;
\ForEach{$i \in [cx_2+1,\; cx_2+2,\; \dots,\; cx_2+n]$}{
    $\mathit{pos} \gets i \bmod n$\;
    \If{$\mathit{pos} \notin [cx_1, cx_2]$}{
        \While{$P_2[j] \in \mathit{seg}$}{
            $j \gets (j + 1) \bmod n$\;
        }
        $C[\mathit{pos}] \gets P_2[j]$\;
        $\mathit{seg}.\text{add}(P_2[j])$\;
        $j \gets (j + 1) \bmod n$\;
    }
}
\Return{$C$}\;
\end{algorithm}

\subsection{Mutazione}

Sono implementati tre operatori di mutazione, scelti casualmente con probabilit\`a
differenziate:
\begin{enumerate}[label=\textbf{\arabic*.}]
    \item \textbf{Swap Mutation} (40\%): Scambia due elementi in posizioni casuali.
          Produce perturbazioni locali, utile per il fine-tuning. \`E complementare
          al 2-opt della ricerca locale (che opera sulle rotte, non sulla permutazione).
    \item \textbf{Insert Mutation} (30\%): Rimuove un elemento e lo reinserisce in
          una posizione diversa. Modifica l'ordine relativo in modo pi\`u drastico
          dello swap, aiutando a ``saltare'' fuori da ottimi locali.
    \item \textbf{Inversion Mutation} (30\%): Inverte un sottosegmento della
          permutazione. Poich\'e il CVRP usa distanze simmetriche, l'inversione
          non altera la fitness della permutazione ma cambia l'output dello Split,
          creando nuove opportunit\`a di partizionamento.
\end{enumerate}

La scelta stocastica dell'operatore evita bias sistematici e mantiene la diversit\`a
delle perturbazioni. Il tasso di mutazione \`e $p_m = 0.236$, quasi doppio rispetto
al default 0.1, riflettendo la necessit\`a di esplorazione aggressiva identificata
da Optuna.

\section{Ricerca Locale Granulare}

L'ibridazione \`e realizzata attraverso quattro operatori di ricerca locale, applicati
con probabilit\`a $p_{ls} = 0.259$ ai nuovi individui. Tutti gli operatori usano
strategia \textbf{steepest descent} (best improvement): ad ogni iterazione viene
valutato l'intero vicinato e viene applicata la mossa che produce il miglioramento
maggiore, anzich\'e accettare il primo miglioramento trovato (first improvement).
Questa scelta privilegia la qualit\`a della soluzione rispetto alla velocit\`a di
esecuzione, ed \`e particolarmente efficace quando combinata con la ricerca locale
granulare che riduce drasticamente la dimensione del vicinato da esplorare.

\begin{enumerate}[label=\textbf{\arabic*.}]
    \item \textbf{2-opt (intra-route)}: Per ogni route, cerca coppie di archi $(a,b)$
          e $(c,d)$ tali che la loro sostituzione con $(a,c)$ e $(b,d)$ riduca il
          costo totale.
    \item \textbf{Or-opt (intra-route)}: Trasferisce segmenti di 1, 2 o 3 nodi
          consecutivi in posizioni diverse all'interno della stessa route.
    \item \textbf{Relocate (inter-route)}: Sposta un cliente da una route all'altra.
    \item \textbf{Exchange (inter-route)}: Scambia due clienti appartenenti a route
          diverse.
\end{enumerate}

\subsection{Educazione a Due Livelli}

Per ottimizzare il rapporto qualit\`a/tempo CPU, l'HGA adotta un'\textbf{educazione
a due livelli}:
\begin{itemize}
    \item \textbf{Light} (applicata al 100\% dei figli): solo 2-opt compilato JIT
          con Numba. Garantisce che ogni figlio abbia almeno le rotte
          ``disincrociate'', eliminando difetti banali a costo quasi nullo.
    \item \textbf{Full} (applicata con probabilit\`a $p_{ls} = 0.259$): tutti e
          quattro gli operatori (2-opt, Or-opt, Relocate, Exchange). \`E la
          cosiddetta ``educazione pesante'', computazionalmente costosa (5 FE per
          chiamata) ma in grado di raffinare significativamente le soluzioni.
\end{itemize}

Questa architettura a due livelli \`e una delle originalit\`a dell'implementazione:
invece di applicare la ricerca locale completa a tutti i figli (con costo proibitivo)
o solo al migliore (con efficacia limitata), si usa il 2-opt come filtro rapido
universale e si riserva la ricerca completa a una frazione casuale dei figli.

\subsection{Ottimizzazioni Computazionali della Ricerca Locale}

L'efficienza della ricerca locale \`e critica per le performance complessive. Sono
state implementate le seguenti ottimizzazioni:
\begin{itemize}
    \item \textbf{Delta cost evaluation} $O(1)$: per valutare se uno spostamento
          \`e vantaggioso, non si ricalcola l'intera soluzione ma solo la differenza
          di costo delle 2 rotte modificate: $\Delta = (\text{costo}_{A, dopo} +
          \text{costo}_{B, dopo}) - (\text{costo}_{A, prima} + \text{costo}_{B,
          prima})$. Questo trasforma un'operazione $O(n)$ in $O(1)$.
    \item \textbf{Pre-calcolo dei carichi}: i load delle route sono mantenuti in
          vettori e aggiornati in $O(1)$ ad ogni mossa, evitando ricalcoli $O(n)$
          a ogni iterazione.
    \item \textbf{Tracciamento solo per indici}: il vicinato viene esplorato
          calcolando solo indici e delta; la mossa fisica (costosa, con copia
          delle route) viene applicata in-place solo alla fine del ciclo, se \`e
          stato trovato un miglioramento.
    \item \textbf{Ordine randomizzato}: Relocate ed Exchange vengono eseguiti in
          ordine casuale a ogni chiamata per evitare bias sistematici.
    \item \textbf{Limite alle iterazioni} ($\text{max\_iter}=2$): limita i cicli
          di ricerca locale a un massimo di 2 passate per chiamata, evitando di
          sprecare budget computazionale su miglioramenti infinitesimali.
\end{itemize}

Queste ottimizzazioni permettono all'algoritmo di raggiungere oltre 6.900 valutazioni
al secondo per core, rendendo fattibili 350.000 FE in circa 50 secondi anche su
istanze da 101 clienti.

Per mitigare l'elevato costo computazionale degli operatori inter-route, \`e stata
integrata la \textbf{Ricerca Locale Granulare} (Granular Local Search, GLS)
\cite{toth2003granular}. Per
ciascun cliente $i$ si definisce un intorno di vicinanza $\Gamma_i$ composto dai soli
$\gamma = 25$ clienti pi\`u vicini. Durante la ricerca locale, un cliente viene
inserito o scambiato solo in posizioni adiacenti a nodi nel suo intorno granulare,
riducendo la complessit\`a a $O(n \cdot \gamma)$.

Senza GLS, la ricerca locale inter-route avrebbe complessit\`a $O(m^2 \cdot n_r^2)$
dove $m$ \`e il numero di rotte ed $n_r$ la lunghezza media. Con GLS, il vicinato si
riduce drasticamente perch\'e \`e improbabile che spostare un cliente in una rotta
geograficamente lontana produca un miglioramento. Il valore $\gamma = 25$, quasi
doppio rispetto al default 15, \`e stato determinato da Optuna: un vicinato pi\`u
ampio paga in termini di qualit\`a senza un costo computazionale proibitivo.

\begin{algorithm}[H]
\footnotesize
\caption{Relocate con Ricerca Locale Granulare (steepest descent)}
\label{alg:relocate}
\KwIn{Insieme di route $\mathcal{R}$, matrice distanze $d$, domande $q$, capacit\`a $Q$, intorni granulari $\Gamma$, depot $0$, max iterazioni $m$}
\KwOut{Route migliorate $\mathcal{R}^*$}

\tcp{Pre-calcolo carichi delle route}
$\mathit{loads}[r] \gets \sum_{n \in r} q[n]$ per ogni route $r \in \mathcal{R}$\;
$\mathit{best} \gets \mathcal{R}$\;
$\mathit{bestCost} \gets \textsc{Cost}(\mathit{best})$\;
$\mathit{iter} \gets 0$\;

\Repeat{$\neg \mathit{improved} \lor (m > 0 \land \mathit{iter} \geq m)$}{
    $\mathit{iter} \gets \mathit{iter} + 1$\;
    $\mathit{improved} \gets \text{false}$\;
    $\mathit{bestMoveCost} \gets \mathit{bestCost}$\;
    $\mathit{bestMove} \gets \text{null}$\;

    \ForEach{route $r_{\text{from}} \in \mathit{best}$}{
        \ForEach{nodo $v$ in posizione $f$ di $r_{\text{from}}$}{
            $\delta_{\text{from}} \gets d[\mathit{prev}(v), \mathit{next}(v)] - (d[\mathit{prev}(v), v] + d[v, \mathit{next}(v)])$\;

            \ForEach{route $r_{\text{to}} \in \mathit{best}$, $r_{\text{to}} \neq r_{\text{from}}$}{
                \If{$\mathit{loads}[r_{\text{to}}] + q[v] > Q$}{
                    \textbf{continue}\;
                }
                \ForEach{posizione $p \in [0, |r_{\text{to}}|]$}{
                    \tcp{Filtro GLS: verifica adiacenza granulare}
                    $\mathit{pn} \gets r_{\text{to}}[p-1]$ se $p > 0$ altrimenti $0$\;
                    $\mathit{nn} \gets r_{\text{to}}[p]$ se $p < |r_{\text{to}}|$ altrimenti $0$\;
                    \If{$(\mathit{pn} \neq 0 \land \mathit{pn} \notin \Gamma[v]) \land (\mathit{nn} \neq 0 \land \mathit{nn} \notin \Gamma[v])$}{
                        \textbf{continue}\;
                    }
                    $\delta_{\text{to}} \gets d[\mathit{pn}, v] + d[v, \mathit{nn}] - d[\mathit{pn}, \mathit{nn}]$\;
                    $\mathit{newCost} \gets \mathit{bestCost} + \delta_{\text{from}} + \delta_{\text{to}}$\;
                    \If{$\mathit{newCost} < \mathit{bestMoveCost}$}{
                        $\mathit{bestMoveCost} \gets \mathit{newCost}$\;
                        $\mathit{bestMove} \gets (r_{\text{from}}, f, r_{\text{to}}, p)$\;
                    }
                }
            }
        }
    }

    \If{$\mathit{bestMove} \neq \text{null}$}{
        $(r_{\text{from}}, f, r_{\text{to}}, p) \gets \mathit{bestMove}$\;
        $v \gets \mathit{best}[r_{\text{from}}][f]$\;
        Rimuovi $v$ da $\mathit{best}[r_{\text{from}}]$ e inseriscilo in $\mathit{best}[r_{\text{to}}]$ alla posizione $p$\;
        $\mathit{loads}[r_{\text{from}}] \gets \mathit{loads}[r_{\text{from}}] - q[v]$\;
        $\mathit{loads}[r_{\text{to}}] \gets \mathit{loads}[r_{\text{to}}] + q[v]$\;
        $\mathit{bestCost} \gets \mathit{bestMoveCost}$\;
        $\mathit{improved} \gets \text{true}$\;
    }
}
\Return{$\mathit{best}$}\;
\end{algorithm}

\section{Elitismo e Rimpiazzo}

I migliori $e = 4$ individui della popolazione corrente sono preservati nella
generazione successiva (elitismo). Il resto della popolazione \`e sostituito
dai nuovi individui generati tramite crossover, mutazione e ricerca locale.

\subsection{Rilevamento Duplicati}

Per mantenere la diversit\`a genetica ed evitare la convergenza prematura, a ogni
generazione viene eseguito il \textbf{rilevamento dei duplicati}: individui con costo
identico (arrotondato alla terza cifra decimale) vengono identificati come cloni e
\textbf{sostituiti con individui casuali freschi}, pre-educati con 2-opt. Senza
questo meccanismo, la popolazione collasserebbe rapidamente su poche soluzioni quasi
identiche (deriva genetica), sprecando il budget computazionale e impedendo ulteriore
miglioramento.

\section{Criterio di Arresto}

L'algoritmo si arresta dopo $FE = 350{,}000$ valutazioni della funzione fitness,
corrispondenti a circa $4{,}300$ generazioni con $\mu = 81$.

\section{Parametri Riepilogativi}

\begin{table}[H]
\centering
\caption{Parametri dell'HGA --- Configurazione Tuned}
\label{tab:params}
\begin{tabular}{@{}ll@{}}
\toprule
\textbf{Parametro} & \textbf{Valore} \\
\midrule
Dimensione popolazione ($\mu$) & 81 \\
Crossover rate ($p_c$) & 0.675 \\
Mutation rate ($p_m$) & 0.236 \\
Local search rate ($p_{ls}$) & 0.259 \\
Tournament size ($k$) & 4 \\
Elite count ($e$) & 4 \\
Granular size ($\gamma$) & 25 \\
Max valutazioni (FE) & 350{,}000 \\
Generazioni medie & $\sim$4{,}300 \\
\bottomrule
\end{tabular}
\end{table}

\section{Ciclo Principale dell'HGA}

L'Algoritmo~\ref{alg:hga} mostra il ciclo evolutivo principale. La popolazione
iniziale viene creata combinando euristiche costruttive e permutazioni casuali.
Ad ogni generazione, i migliori $e$ individui sono preservati (elitismo), mentre
i restanti sono generati tramite selezione a torneo, Order Crossover (OX), mutazione
e ricerca locale. L'algoritmo termina dopo $FE = 350{,}000$ valutazioni della
funzione fitness.

\begin{algorithm}[H]
\footnotesize
\caption{Ciclo Principale dell'HGA}
\label{alg:hga}
\KwIn{Istanza CVRP; parametri: $\mu$, $p_c$, $p_m$, $p_{ls}$, $k$, $e$, $\gamma$, $FE$}
\KwOut{Migliore soluzione $S^*$}

\tcp{Inizializzazione}
$\mathcal{P} \gets \emptyset$\;
$\mathcal{P}.\text{append}(\textsc{Split}(\textsc{NearestNeighbor}()))$\;
$\mathcal{P}.\text{append}(\textsc{Split}(\textsc{ClarkeWrightSavings}()))$\;
\While{$|\mathcal{P}| < \mu$}{
    $\mathcal{P}.\text{append}(\textsc{Split}(\textsc{RandomPermutation}()))$\;
}
$S^* \gets \arg\min_{S \in \mathcal{P}} S.\text{cost}$\;
$\mathit{evals} \gets |\mathcal{P}|$\;

\tcp{Ciclo evolutivo}
\While{$\mathit{evals} < FE$}{
    $\mathcal{P}' \gets \text{top-}e(\mathcal{P})$\;  \tcp{Elitismo}
    \While{$|\mathcal{P}'| < \mu$}{
        \tcp{Selezione e Crossover}
        \If{$\text{random}() < p_c$}{
            $P_1 \gets \textsc{TournamentSelect}(\mathcal{P}, k)$\;
            $P_2 \gets \textsc{TournamentSelect}(\mathcal{P}, k)$\;
            $\pi_{\text{child}} \gets \textsc{OrderCrossover}(P_1, P_2)$\;
        }
        \Else{
            $P \gets \textsc{TournamentSelect}(\mathcal{P}, k)$\;
            $\pi_{\text{child}} \gets \text{flatten}(P.\text{routes})$\;
        }

        \tcp{Mutazione}
        \If{$\text{random}() < p_m$}{
            $\pi_{\text{child}} \gets \textsc{Mutate}(\pi_{\text{child}})$\;  \tcp{Swap 40\%, Insert 30\%, Inversion 30\%}
        }

        \tcp{Decodifica}
        $S_{\text{child}} \gets \textsc{Split}(\pi_{\text{child}})$\;

        \tcp{Educazione leggera (sempre)}
        $S_{\text{child}} \gets \textsc{TwoOpt}(S_{\text{child}})$\;

        \tcp{Ricerca locale completa (probabilistica)}
        \If{$\text{random}() < p_{ls}$}{
            $S_{\text{child}} \gets \textsc{TwoOpt}(S_{\text{child}})$\;
            $S_{\text{child}} \gets \textsc{OrOpt}(S_{\text{child}})$\;
            $S_{\text{child}} \gets \textsc{Relocate}(S_{\text{child}}, \gamma)$\;
            $S_{\text{child}} \gets \textsc{Exchange}(S_{\text{child}}, \gamma)$\;
        }

        $\mathcal{P}'.\text{append}(S_{\text{child}})$\;

        \If{$S_{\text{child}}.\text{cost} < S^*.\text{cost}$}{
            $S^* \gets S_{\text{child}}$\;
        }
    }

    \tcp{Rimpiazzo duplicati}
    \ForEach{coppia di individui con stesso costo in $\mathcal{P}'$}{
        Sostituisci uno con $\textsc{Split}(\textsc{RandomPermutation}())$\;
    }

    $\mathcal{P} \gets \mathcal{P}'$\;
}

\Return{$S^*$}\;
\end{algorithm}

\section{Architettura Software}

L'algoritmo \`e integrato in un'architettura moderna con backend Python (FastAPI)\cite{fastapi} e
frontend React con visualizzazione in tempo reale tramite WebSocket. Il backend espone
endpoint REST per avviare l'ottimizzazione e recuperare i risultati, mentre il frontend
mostra simultaneamente la disposizione geografica delle route, i grafici di convergenza
e le statistiche aggregate. La parallelizzazione multi-run sfrutta \texttt{ThreadPoolExecutor}
per eseguire run indipendenti in parallelo, e l'accelerazione Numba JIT \cite{lam2015numba} per i calcoli
critici (matrici di distanza, Split DP).
